% ==================================================
% CHAPTER 4: IMPLEMENTATION
% ==================================================

\chapter{Implementation}
\label{ch:implementation}

\section{Code Base}

The backend is a Python package, \texttt{hostelswap}, installed from
\texttt{code/} with \texttt{pyproject.toml}. Table~\ref{tab:modules}
lists its modules.

\begin{table}[H]
\centering
\small
\begin{tabularx}{\textwidth}{>{\ttfamily\raggedright\arraybackslash}p{0.28\textwidth} r L}
\toprule
\textrm{\textbf{Module}} & \textbf{Lines} & \textbf{Responsibility} \\
\midrule
domain/models.py, preferences.py & 89 & Rooms, slots, allocations; criteria and weighted hard/soft preferences \\
domain/pool.py & 66 & The closed pool of students, rooms and slots in one round \\
domain/scoring.py & 65 & How well a room meets a student's preferences \\
domain/matrix.py & 70 & Students-by-slots cost matrix \\
domain/solver.py & 73 & Hungarian assignment (SciPy) \\
domain/kbest.py & 65 & Murty's $k$ best assignments \\
domain/cycles.py & 53 & Splitting an assignment into swap chains \\
domain/feasibility.py & 45 & Pareto filter: nobody worse off, someone better off \\
domain/ranking.py & 238 & Options that differ in kind; chains offered per student; chain cap \\
domain/pipeline.py & 33 & Full run: pool in, options or chains out \\
domain/proposal.py & 127 & Approval state machine \\
synth/generator.py & 199 & Seeded synthetic cohorts for tests and benchmarks \\
rounds.py & 238 & Round workflow: enrol, prioritise, lock and run, offer, respond, expire \\
db.py & 1193 & All SQL, connection pool, settings \\
auth.py & 95 & Password hashing and session cookies \\
api.py & 613 & FastAPI routes and role checks \\
demo\_data.py & 155 & Fixed 30-student demo cohort \\
\bottomrule
\end{tabularx}
\caption{Backend modules}
\label{tab:modules}
\end{table}

Around the package sit \texttt{db/} (migrations, the migration runner and
seed scripts), \texttt{frontend/} (the React application),
\texttt{tests/} (152 tests), \texttt{benchmarks/} and \texttt{docs/}
(the documentation site and journals).

\section{Matching Engine}

Figure~\ref{fig:pipeline} shows the matching pipeline. Every stage is a
pure function of the previous stage's output.

\begin{figure}[H]
\centering
\begin{tikzpicture}[node distance=4mm]
  \node[stepbox] (s1) {Pool: students, slots, preferences};
  \node[stepbox, below=of s1] (s2) {Cost matrix (scores, $\infty$ if infeasible)};
  \node[stepbox, below=of s2] (s3) {Hungarian: best assignment};
  \node[stepbox, below=of s3] (s4) {Murty: next $K-1$ assignments};
  \node[stepbox, below=of s4] (s5) {Cycle decomposition: chains};
  \node[stepbox, below=of s5] (s6) {Pareto filter, chain cap};
  \node[stepbox, below=of s6] (s7) {Up to 3 chains per student};
  \foreach \a/\b in {s1/s2,s2/s3,s3/s4,s4/s5,s5/s6,s6/s7} \draw[flow] (\a) -- (\b);
\end{tikzpicture}
\caption{Matching pipeline}
\label{fig:pipeline}
\end{figure}

\subsection{Preferences and Scoring}

A preference names a criterion (room, hostel, floor, direction,
washroom, room type or roommate), a value, a weight and whether it is
hard. The score of a room for a student is the share of their soft
preference weight that the room meets:
\[
  \mathrm{score}(s, r) = \frac{\sum_{p \in \mathrm{soft}(s),\ p \text{ met by } r} w_p}
                              {\sum_{p \in \mathrm{soft}(s)} w_p}
\]
An unmet \emph{hard} preference does not lower the score: it makes the
room infeasible for that student. A student with no soft preferences is
content anywhere (score 1). The score keeps the lists of met and unmet
preferences, so every option can explain itself in the interface
(``floor 3: met; attached washroom: missed'').

\subsection{Cost Matrix and the Hungarian Algorithm}

Rows are students and columns are bed slots. Each cell is the negated
score, because the solver minimises, or $\infty$ if a hard preference is
unmet. A tiny inertia term ($10^{-9}$) is added to every slot other than
the student's own. Without it, the solver would happily move a student
to the other bed of the same room, which changes nothing for them and
shows up as a meaningless swap.

The matrix is solved with
\texttt{scipy.optimize.linear\_sum\_assignment}, an implementation of
the Hungarian method \cite{kuhn1955hungarian,munkres1957}, in $O(n^3)$.
If hard preferences cannot all be met at once, the solver reports that
no feasible assignment exists.

\subsection{Murty's Algorithm}

The best assignment is often not the best answer. It may demote one
student to help several, which the Pareto filter then rejects, or it may
need one enormous chain. Murty's algorithm \cite{murty1968ranking}
enumerates the next best assignments. Any assignment other than the best
must exclude at least one of the best one's edges, so for each edge in
turn the algorithm forbids it, requires the edges before it, and
re-solves. These subproblems are disjoint and together cover every
other assignment, so nothing repeats and nothing is missed. A priority
queue returns the assignments cheapest first. A live round uses
$K = 50$.

Each returned assignment creates up to $n$ subproblems, and each is a
Hungarian solve, so $K$ results cost $O(K \cdot n \cdot n^3) = O(Kn^4)$
in the worst case. The subproblems differ from their parent by fixed
(``required'') edges. The first version forced a fixed edge by setting
the rest of its row and column to $\infty$, so every subproblem was
still a full $n \times n$ solve, and it found each student's row with
an $O(n)$ search of a tuple. The solver now deletes the fixed rows and
columns, so a subproblem with $r$ fixed edges is an
$(n-r) \times (n-r)$ solve, and the row and column of each student and
slot are looked up in dictionaries built once per matrix. The results
are identical; Section~\ref{sec:performance} gives the speed-up.

\subsection{Chains and the Pareto Filter}

Because the pool is closed, each assignment is a permutation of students
over the slots they already hold, and every permutation splits uniquely
into disjoint cycles. Each cycle of length two or more is a swap chain:
student 1 moves into student 2's slot, student 2 into student 3's, and
the last into student 1's.

A chain is kept only if it is a \textbf{Pareto improvement}: every
member's score in the new slot is at least their score in the current
one, and at least one member's is strictly higher. The check runs after
the solve, not inside it, because the solver maximises the total, which
can legitimately mean demoting one student to help several. Such a
chain would only be rejected at the approval step, so it is never
offered.

\subsection{Options per Student}

The prototype picked one plan for the whole hostel: the assignment with
the highest cohort satisfaction. That answer hands most students
nothing. A two-person swap in which both reach 100\% can beat a rotation
that lifts twenty. The final engine answers a different question:
\emph{what is the best move available to me?}

\texttt{candidate\_chains} collects every distinct Pareto-improving chain
from all $K$ assignments, drops chains longer than the cap, and keeps,
for each student, their own three best: ordered by how well the chain
serves that student, then by fewer members. The chains shown to
different students may overlap and contradict each other on purpose.
Whoever offers first takes the rooms; the unique index on live
proposals stops a second chain from claiming the same students, and the
other students fall back to their next option.

The original whole-hostel ranking, which picks up to three options that
\emph{differ in kind} (best overall match, fastest to execute, best
location fit), is kept in \texttt{rank}. It awards a label only to an
option that actually beats the others on that measure, so it may return
fewer than three rather than a misleading label.

\subsection{Chain-Length Cap}

Every member of a chain has to accept before anyone moves, so a chain of
fourteen is one offer with fourteen chances to fall through. The engine
accepts an optional \texttt{max\_chain}; rounds use eight by default,
and a warden can override it when locking a round. The value was chosen
by benchmark (Section~\ref{sec:chain-cap}). A long cycle cannot be split
into shorter ones, so the chains a cap removes are lost, not
redistributed.

\section{Round Workflow}

\texttt{rounds.py} drives a round through its lifecycle and holds the
rules about which action is allowed in which state:

\begin{enumerate}
  \item \textbf{Open}: only a \texttt{draft} round can be opened.
  \item \textbf{Enrol and prioritise}: only while the round is
  \texttt{open}. Priorities are saved as a new active preference set.
  \item \textbf{Lock and run}: the round becomes \texttt{locked} and
  then \texttt{running}; the engine runs on the round's pool; the chains
  are stored in \texttt{round\_chain\_options}; the round becomes
  \texttt{completed}.
  \item \textbf{Offer}: a member of a stored chain turns it into a
  proposal with a deadline three days ahead. Before inserting it, the
  workflow sweeps every overdue proposal to \texttt{expired}, so a chain
  nobody answered cannot block this one.
  \item \textbf{Respond}: the workflow first checks the deadline. A late
  reply marks the proposal expired and is refused with ``this proposal
  has expired''. Otherwise the reply goes through the domain state
  machine: a rejection settles the proposal and frees every member; the
  last approval executes it and appends the new allocations.
  \item \textbf{Cancel}: a warden can cancel a draft, open or completed
  round, which withdraws its open proposals and frees their members in
  one transaction, or withdraw a single ring. Executed swaps stay.
\end{enumerate}

The student page counts down each offer's deadline and shows an overdue
offer as expired at once, without waiting for the server's next sweep.

Mistakes by the client, such as a double click or a stale page that
answers an already-settled proposal, raise \texttt{InvalidTransition} in
the domain and are returned as \texttt{409 Conflict}, not as server
errors.

\section{Persistence}

\texttt{db.py} uses psycopg~3 with a connection pool, so a request does
not pay for a new TLS connection to Supabase. The student's round view
is loaded with one query instead of one per section. Settings come from
the environment first and then from a \texttt{.env} file at the
repository root.

The schema lives in \texttt{db/migrations/} as numbered SQL files.
\texttt{db/migrate.py} applies every pending file in order, each in its
own transaction, and records it in \texttt{schema\_migrations}. An
applied file is never edited; a change is a new file. A
\texttt{--baseline} option marks the files that created an existing
database by hand as applied, without running them again.

Two seed scripts produce realistic data: \texttt{seed\_hostels.py} loads
four hostels and 1000 students (it refuses to clear the database without
\texttt{--reset}), and \texttt{seed\_round.py} adds an enrolled cohort
with priorities to the open round, skipping roll numbers that already
exist, so it can be run again safely.

\section{Web Application}

The React application (\texttt{frontend/src}) has a typed API client
(\texttt{api.ts}) and shared types (\texttt{types.ts}) that mirror the
JSON the backend returns. The student page shows each option as a chain
card: who moves where, the student's match percentage before and after,
and the preferences met and missed. Offering a chain, setting priorities
and choosing room types happen in dialogs, so the user never leaves the
page. Roommate requests are sent by roll number and must be accepted by
the other student before the engine uses them.

\section{Documentation}

The repository contains a documentation site built with Material for
MkDocs and published to GitHub Pages by a workflow on every push to
\texttt{main}. It holds the project overview, setup instructions and the
weekly journals of each member, written as problem, cause and fix.
